% !PRINTED 78
\weberpage{78}

Alle primären Theiler des vollen Kreistheilungskörpers $\Omega_m$ sind demnach in der Form $\Omega(\eta)$ darstellbar, d.\,h. alle Zahlen eines solchen Theilers sind rational durch die Kreistheilungsperioden $\eta$ darstellbar, und da man die nicht primären Theiler von $\Omega_m$ als primäre Theiler von niedrigeren Kreistheilungskörpern wiederfindet, so folgt:

\gesperrt{4. Alle Kreistheilungskörper sind in der Form $\Omega(\eta)$ darstellbar.}

Oder auch:

\gesperrt{5. Jede rationale Function von Einheitswurzeln kann als rationale Function einer Kreistheilungsperiode dargestellt werden.}

Und dieser Satz lässt sich auch in der Form aussprechen:

\gesperrt{6. Ist $\mathfrak{A}$ ein beliebiger primärer oder nicht primärer Theiler von $\mathfrak{N}$ vom Index $e$, so giebt es einen Theiler $m_1$ von $m$ und eine in $\Omega_{m_1}$ gelegene Kreistheilungsperiode $\eta$ vom Index $e$, so dass $\eta$, als Function von $r$ aufgefasst, eine zur Gruppe $\mathfrak{A}$ gehörige Function ist, also, wenn $a$ zu $\mathfrak{A}$ gehört, durch die Substitution $(r, r^a)$ ungeändert bleibt und für alle Substitutionen von $\mathfrak{N}$ $e$ verschiedene Werthe erhält.}

Durchläuft $a$ eine Gruppe $\mathfrak{A}$, die ein primärer Theiler von $\mathfrak{N}$ ist, und setzen wir

\[
\eta_h = \sum^a r^{ah},
\]

auch wenn $h$ nicht relativ prim zu $m$ ist, so können wir diese Grössen $\eta_h$, die alle die Permutationen der Gruppe $\mathfrak{A}$ gestatten, rational durch $\eta = \sum r^a$ darstellen. Andererseits lässt sich aber auch jede rationale Function von $\eta$ linear durch die $\eta_h$ darstellen. Dies wird bewiesen sein, wenn gezeigt ist, dass man das Product zweier $\eta_h$ linear durch die $\eta_h$ ausdrücken kann. Es ist aber:

\[
\eta_h \eta_k = \sum^{a, a'} r^{ha + ka'};
\]

nimmt man zuerst die Summe nach $a'$ bei feststehendem $a$, so kann man $a'$ durch $aa'$ ersetzen, da $aa'$ zugleich mit $a'$ die Gruppe $\mathfrak{A}$ durchläuft. Man erhält so

\[
\eta_h \eta_k = \sum^{a, a'} r^{a(h + ka')} = \sum^{a'} \eta_{h + ka'}.
\]
